\section{Full ND-LAr and ND 2\texorpdfstring{$\times$}{x}2}
\label{sec:nd2x2}

\subsection{Two detector sizes, two channel counts}

The full ND-LAr design has 35 optically separated modules, each with two
cathode-separated drift regions. Each module contains 60 LCM panels read
out by two SiPM channels each and 20 ArCLight panels read out by six each
\cite{nd_cdr}. Thus
\begin{equation}
\boxed{N_{\mathrm{ch}}^{\mathrm{ND-LAr}}
=35(60\times2+20\times6)=8400,}
\end{equation}
with 4200 channels of each technology. This also matches the explicitly
quoted 8400 electronic channels in the 2024 design presentation
\cite{nd_channels}. The count refers to PDS readout channels, not charge
pixels, individual photon hits, or the other near-detector subsystems.
The older CDR computing estimate (section 9.3.2.1) instead assumes about
4200 photon-detector channels; this comparison uses the explicitly quoted
8400-channel electronics design, consistent with reading the individual
SiPM signals described in its light-readout section.

The 2$\times$2 demonstrator has four smaller modules, eight TPC regions,
and 384 individually digitized SiPM channels~\cite{nd_operation}.
The alternating three-LCM/one-ArCLight arrangement contains 96 LCMs
and 32 ArCLights in total, giving
\begin{equation}
\boxed{N_{\mathrm{ch}}^{2\times2}
=96\times2+32\times6=384,}
\end{equation}
with 192 channels of each technology. A full ND-LAr module has 240 light
channels; a smaller prototype module has 96. Scaling 384 by 35/4 would
therefore give the wrong full-detector count.

\subsection{Argon activity and optical compartments}

For full ND-LAr, the CDR active-volume reference is
$35\times(1\times1\times3)=\SI{105}{m^3}$~\cite{nd_cdr}. At the adopted
 density this corresponds to \SI{146.475}{t} and
$R_{39,\mathrm{decay}}=\SI{1.412019e5}{s^{-1}}$.
For 70 equal optically isolated TPC regions, each has \SI{1.5}{m^3} of
active argon and about 2017 decays/s. Use active, optically emitting mass
here, rather than the 50--67 t fiducial analysis masses quoted for particular
selections. Optically visible external bath argon is excluded.

For 2$\times$2, use the \SI{2.4}{t} active mass and 29\% geometrical
coverage reported in the operation paper~\cite{nd_operation}, updating the
older 2.6 t/25\% inputs. Its total activity is
$(0.964)(2400)=\SI{2313.6}{s^{-1}}$, or \SI{289.2}{s^{-1}} per equal TPC.
The active volume per region is then
$(2400/8)/1395=\SI{0.215}{m^3}$; external module envelopes must not be
 treated as fully active while retaining this mass normalization.

\subsection{From total photons to each technology's channel rate}

For a uniform source in a TPC, the exact direct-light expression is
\begin{equation}
R_{39,j}=S_{39}M_{\mathrm{TPC}}\overline N_{\gamma,39}\epsilon_j
\mathcal G_j,\qquad
\mathcal G_j=\frac{1}{V_{\mathrm{TPC}}}\int_{V_{\mathrm{TPC}}}
\frac{\Omega_j(\bm r)}{4\pi}\,dV.
\label{eq:ND39exact}
\end{equation}
For a rectangular panel at positive normal distance $d$, write signed corner
coordinates relative to the source projection as $X_\pm,Y_\pm$. Then
\begin{equation}
\Omega_j=\sum_{p=\pm1}\sum_{q=\pm1}pq\,
\tan^{-1}\!\left[\frac{X_pY_q}
{d\sqrt{d^2+X_p^2+Y_q^2}}\right].
\label{eq:NDsolidangle}
\end{equation}
For transport beyond unobstructed direct light, use the full response kernel
in Eq.~\eqref{eq:master}. A measured panel PDE is a whole-panel response;
 do not apply it independently to every SiPM and multiply the detected light.

The order-of-magnitude table uses a simpler \emph{uniform first-hit wall-flux}
model: $h_t$ is the fraction of the total interior surface instrumented by
technology $t$, treated as its photon interception probability. No reflected
second encounters are included. This gives technology-average rates
\begin{equation}
\boxed{R_{39,\mathrm{ch},t}\simeq
\frac{S_{39}M\overline N_{\gamma,39}h_t\epsilon_t}
{N_{\mathrm{ch},t}},\qquad t=A,L.}
\label{eq:NDchannel}
\end{equation}
Surface coverage need not equal interception probability for real wall flux;
 the exact geometry integral above replaces this approximation.

For the CDR full-size panels, ArCLight is approximately $0.30\times0.50$ m
and LCM $0.10\times0.50$ m~\cite{nd_cdr}. Their areas per module are
$20(0.15)=3.0$ m$^2$ and $60(0.05)=3.0$ m$^2$. Two isolated
$0.5\times1\times3$ m boxes have total interior area
$2\times2(0.5\times1+0.5\times3+1\times3)=20$ m$^2$,
counting both faces of the cathode. Hence $h_A=h_L=0.15$ and total
coverage is 0.30 in this reference geometry.

For 2$\times$2, the operation paper quotes 29\% coverage,
$0.28\times0.30$ m ArCLights and $0.10\times0.30$ m LCMs
\cite{nd_operation}. The instrumented-area shares are
$f_A=32(0.084)/[32(0.084)+96(0.030)]=14/29$ and $f_L=15/29$.
Therefore $h_A=0.29f_A=0.14$ and $h_L=0.15$.
With $\epsilon_A=0.002$, $\epsilon_L=0.006$~\cite{nd_status} and the
4400-photon reference yield, Eq.~\eqref{eq:NDchannel} gives
\begin{center}
\begin{tabular}{lrrr}
\toprule
Detector & ArCLight & LCM & All-channel mean \\
 & PE/s/channel & PE/s/channel & PE/s/channel \\
\midrule
Full ND-LAr & 44.4 & 133.1 & 88.8 \\
ND 2$\times$2 & 14.8 & 47.7 & 31.3 \\
\bottomrule
\end{tabular}
\end{center}
The corresponding detector totals are $7.46\times10^5$ and
$1.20\times10^4$ primary PE/s. The mean is weighted by electronic
channel counts, not by panel counts. Rescale each technology row with its
actual yield, PDE, and interception probability; the full-size PDE is a
reference extrapolation from prototype values, not a full-size calibration.

The average full-ND channel has about 2.84 times the radioactive-light rate
of a 2$\times$2 channel in this model, while the channel count is 21.875
 times larger. That difference follows from source mass, coverage, and
channel sharing, rather than from simply dividing or multiplying by the
number of modules. The separate site-dependent cosmic estimates are
 given in Section~\ref{sec:cosmicnumbers}.
